\subsection{Two-Stage Miller Op-Amp}\label{sec:amp_two_stage}

\subsubsection{Overview}
A 5T OTA (first stage, transconductance \(g_{m1}\), output conductance
\(G_1 = g_{o2} + g_{o4}\), node capacitance \(C_1\)) followed by a common
source stage \(M_6\) with current source load \(M_7\) (\(G_2 = g_{o6} + g_{o7}\),
load \(C_L\)). The gain is the product of two stages; a Miller capacitor
\(C_c\) across the second stage makes the first pole dominant
(compensation). Its small signal model is the CS stage of
Section~\ref{sec:amp_cs} with \(C_c\) in the place of \(C_{gd}\).

\subsubsection{Schematic}
\begin{circuitfig}
  \draw (0,0) node[draw, minimum width=1.6cm, minimum height=1.2cm](S1){$g_{m1}$, $G_1$};
  \draw (S1.west) ++(0,0.3) -- ++(-0.8,0) node[circ]{} node[left]{$V_{i1}$};
  \draw (S1.west) ++(0,-0.3) -- ++(-0.8,0) node[circ]{} node[left]{$V_{i2}$};
  \draw (S1.east) -- (2.4,0) node[circ]{} node[below left]{$V_1$} to[capacitor, l_={$C_1$}] (2.4,-2) node[ground]{};
  \draw (5,0) node[nmos](M6){}; \node[right] at (5.3,0) {$M_6$};
  \draw (2.4,0) -- (M6.G);
  \draw (M6.S) -- (5,-2) node[ground]{};
  \draw (M6.D) -- (5,1.4) node[circ]{} to[american current source, l_={$I_7$, $G_{o7}$}] (5,3.4) node[vcc]{$V_{dd}$};
  \draw (2.4,0) -- (2.4,1.4) to[capacitor, l={$C_c$}] (5,1.4);
  \draw (5,1.4) -- (7,1.4) node[circ]{} node[above]{$V_{out}$} to[capacitor, l={$C_L$}] (7,-2) node[ground]{};
\end{circuitfig}

\subsubsection{Transfer Function}
\begin{enumerate}
  \item \textbf{Large Signal, Time Domain, DC Input Analysis (Q5)}:
    for zero systematic offset the first stage output \(V_1 = V_{dd} - |V_{GS3}|\)
    must equal the \(V_{GS6}\) that makes \(M_6\) carry \(I_7\). With
    \(|V_{GS3}| = |V_{GS4}|\) this requires
    \[
      \frac{(W/L)_6}{(W/L)_4} = 2\frac{(W/L)_7}{(W/L)_5}
    \]
    (\(M_6\) sees the same \(V_{GS}\) relation as \(M_4\)).
  \item \textbf{Small Signal, Frequency Domain, AC Input Analysis (Q4)}:
    the first stage injects \(g_{m1} V_{id}\) into \(V_1\).\\
    Applying KCL at \(V_1\):
    \begin{align*}
      V_1 (G_1 + sC_1) + (V_1 - V_{out}) sC_c &= g_{m1} V_{id} \\
      V_1 (G_1 + s(C_1 + C_c)) - V_{out} sC_c &= g_{m1} V_{id} \tag{1}
    \end{align*}
    Applying KCL at \(V_{out}\):
    \begin{align*}
      V_{out}(G_2 + sC_L) + (V_{out} - V_1) sC_c + g_{m6} V_1 &= 0 \\
      V_1 &= -\frac{G_2 + s(C_L + C_c)}{g_{m6} - sC_c} V_{out} \tag{2}
    \end{align*}
    Substituting (2) into (1):
    \begin{align*}
      -V_{out}\left(\frac{(G_2 + s(C_L + C_c))(G_1 + s(C_1 + C_c)) + sC_c (g_{m6} - sC_c)}{g_{m6} - sC_c}\right) &= g_{m1} V_{id} \\
      A(s) = \frac{V_{out}}{V_{id}} &= \frac{-g_{m1}(g_{m6} - sC_c)}{G_1 G_2 + s\left[G_2(C_1 + C_c) + G_1(C_L + C_c) + g_{m6} C_c\right] + s^2\left[C_1 C_L + C_1 C_c + C_L C_c\right]}
    \end{align*}
    (the sign is that of the inverting second stage; swapping \(V_{i1}\),
    \(V_{i2}\) restores a positive gain).
\end{enumerate}

\subsubsection{Analysis}
\begin{enumerate}
  \item \textbf{DC Gain}: \(A_{V0} = \frac{g_{m1} g_{m6}}{G_1 G_2} = \frac{g_{m1}}{g_{o2} + g_{o4}} \cdot \frac{g_{m6}}{g_{o6} + g_{o7}}\).
  \item \textbf{Poles} (widely spaced, \(g_{m6} C_c\) dominating \(a_1\)):
    \begin{align*}
      p_1 &\approx \frac{G_1 G_2}{g_{m6} C_c} \hspace{1cm} [C_c \text{ Miller multiplied by } \tfrac{g_{m6}}{G_2}] \\
      p_2 &\approx \frac{g_{m6} C_c}{C_1 C_L + C_1 C_c + C_L C_c} \xrightarrow{C_c \gg C_1} \frac{g_{m6}}{C_L + C_1} \approx \frac{g_{m6}}{C_L}
    \end{align*}
  \item \textbf{Zero}: \(z = +\frac{g_{m6}}{C_c}\) (right half plane, adds phase lag).
  \item \textbf{Gain-Bandwidth}: \(2\pi GB = A_{V0}\, p_1 = \frac{g_{m1} g_{m6}}{G_1 G_2} \cdot \frac{G_1 G_2}{g_{m6} C_c} = \frac{g_{m1}}{C_c}\).
  \item \textbf{Compensation}: in unity feedback (\(F = 1\)) with \(p_1 \to 0\),
    \(z \to \infty\): \(A(s) \approx \frac{2\pi GB \cdot p_2}{s(s + p_2)}\),
    \begin{align*}
      \frac{A}{1 + A} &= \frac{2\pi GB\, p_2}{s^2 + p_2 s + 2\pi GB\, p_2} \hspace{1cm}
      \omega_n = \sqrt{2\pi GB\, p_2} \hspace{1cm} Q = \sqrt{\frac{2\pi GB}{p_2}}
    \end{align*}
    \(Q = 0.707\) (no peaking) needs \(p_2 = 2 \cdot 2\pi GB\);
    phase margin \(PM = 90^\circ - \tan^{-1}\frac{2\pi GB}{p_2} - \tan^{-1}\frac{2\pi GB}{z}\)
    gives \(PM = 63.4^\circ\) at \(p_2 = 2 \cdot 2\pi GB\), \(z \to \infty\).
    With \(z = 10 \cdot 2\pi GB\) (\(g_{m6} = 10 g_{m1}\)): \(p_2 = 2.2 \cdot 2\pi GB\)
    needs \(C_c \geq 0.22 C_L\) for \(PM \approx 60^\circ\).
  \item \textbf{Nulling Resistor}: \(R_z\) in series with \(C_c\) replaces
    \(sC_c\) by \(\frac{sC_c}{1 + sR_z C_c}\); the zero moves to
    \[
      z = \frac{1}{C_c\left(\frac{1}{g_{m6}} - R_z\right)}
    \]
    \(R_z = \frac{1}{g_{m6}}\) removes it, \(R_z > \frac{1}{g_{m6}}\) puts it in the
    left half plane where it can cancel \(p_2\).
  \item \textbf{Slew Rate (Q6)}: \(SR = \min\left(\frac{I_5}{C_c}, \frac{I_7 - I_5}{C_L}\right)\).
\end{enumerate}

\subsubsection{Question}
\(C_L = 5\,\)pF, \(GB = 10\,\)MHz, \(PM = 60^\circ\), \(g_{m6} = 10 g_{m1}\).
Find \(C_c\), \(g_{m1}\), \(g_{m6}\) and check \(p_2\).
\paragraph{Solution:}
\begin{align*}
  C_c &= 0.22 \cdot 5\,pF = 1.1\,pF \hspace{1cm}
  g_{m1} = 2\pi \cdot 10\,MHz \cdot 1.1\,pF = 69.1\,\mu S \hspace{1cm}
  g_{m6} = 691\,\mu S \\
  \frac{p_2}{2\pi} &= \frac{691\,\mu S}{2\pi \cdot 5\,pF} = 22\,MHz = 2.2\,GB \hspace{1cm}
  \frac{z}{2\pi} = \frac{691\,\mu S}{2\pi \cdot 1.1\,pF} = 100\,MHz \\
  PM &= 90^\circ - \tan^{-1}\frac{1}{2.2} - \tan^{-1}\frac{1}{10} = 90 - 24.4 - 5.7 = 59.9^\circ
\end{align*}
